\chapter*{Resumo}
\addcontentsline{toc}{chapter}{Resumo}

A priorização de requisitos de software é uma das atividades mais
sensíveis da Engenharia de Requisitos: técnicas tradicionais como a
\textbf{Votação Acumulativa} e o \textbf{PERT/CPM} --- descritas em
detalhe no Capítulo~\ref{ch:fundamentos} --- dependem fortemente da
perceção subjetiva dos \textit{stakeholders} e tendem a tornar-se
inconsistentes em sistemas com elevado número de requisitos
interdependentes \parencite{karlsson2002, silva2018}. Esta dissertação
investiga uma alternativa de natureza estrutural, assente na aplicação
de algoritmos de ranqueamento de grafos --- PageRank, HITS e SALSA ---
ao grafo de dependências entre requisitos.

Foi desenvolvido o \textbf{ReqGraph}, um protótipo \emph{open source}
em Python que automatiza o \textit{pipeline}
\emph{Código-fonte \(\rightarrow\) AST \(\rightarrow\) Call Graph
\(\rightarrow\) Grafo de Requisitos}, e o módulo \code{ranker.py}, que
aplica os três algoritmos ao grafo extraído. A validação foi conduzida
sobre três projetos reais de código aberto --- \emph{CPython stdlib},
\emph{Flask} e \emph{scikit-learn} --- selecionados por
representarem regimes crescentes de acoplamento arquitetural. A
comparação entre os rankings produzidos pelos três algoritmos é feita
de forma quantitativa, recorrendo aos coeficientes de
\textbf{correlação de Spearman} (\(\rho\)) e \textbf{Kendall}
(\(\tau_b\)). Os resultados mostram que, em grafos com densidade
suficiente, os algoritmos partilham parcialmente a identificação dos
requisitos mais centrais; em grafos esparsos, divergem de forma
informativa, oferecendo perspetivas complementares de importância
--- transitiva (PageRank), direta (HITS \emph{authority}) e de
orquestração (HITS \emph{hub}). Conclui-se que o ranqueamento estrutural
constitui um complemento robusto às técnicas tradicionais, contribuindo
para a redução --- não a eliminação --- do viés subjetivo na priorização
de requisitos. Discutem-se em detalhe as limitações da abordagem,
nomeadamente a subjetividade residual do mapeamento
\code{func\_to\_req} e a dimensão reduzida do corpus.

\vspace{0.3cm}
\noindent\textbf{Palavras-chave:} Engenharia de Requisitos; Priorização
Estrutural; Análise Estática de Código; Grafos; PageRank; HITS; SALSA;
Correlação de Spearman.
